% GENERATED FILE. Edit canonical lessons, then run npm run generate:books.
% canonical-source-hash: fnv1a64:f062e807b67bc43b
% canonical-lessons: ML-C449-niyantrikkuka, ML-C449-varddhikkuka, ML-C449-varddhippikkuka, ML-C449-uyaruka, ML-C449-uyarttuka

\chapter{Up and Down}
\label{ch:ml-up-and-down}
\hlchaptermodality{449}

\begin{chapteropening}
\textbf{By the end of this chapter:} \emph{I can say to control, to go up, to increase, to rise and to raise.}

Complete the last lesson of chapter 449: I can say to control, to go up, to increase, to rise and to raise.
\end{chapteropening}

\section[\texorpdfstring{\ml{നിയന്ത്രിക്കുക} (niyantrikkuka)}{നിയന്ത്രിക്കുക (niyantrikkuka)}]{\ml{നിയന്ത്രിക്കുക} (niyantrikkuka) — to control}

\label{lesson:ML-C449-niyantrikkuka}



\begin{quote}
Before the new one: say the Malayalam for to analyse, then the Malayalam for a method.
\end{quote}

\subsection*{You'll want to know: \ml{നിയന്ത്രിക്കുക}}
\textbf{\ml{നിയന്ത്രിക്കുക}} — \emph{niyantrikkuka} — \textquotedblleft{}to control\textquotedblright{}.

\begin{quote}
\ml{പോലീസ്} \ml{ആൾക്കൂട്ടത്തെ} \ml{നിയന്ത്രിച്ചു}.
\end{quote}

\begin{quote}
\emph{pōlīsŭ āḷkkūṭṭatte niyantriccu.}
\end{quote}

\begin{quote}
The police controlled the crowd.
\end{quote}

Past: \ml{നിയന്ത്രിച്ചു}, \emph{niyantriccu}.

The noun is \ml{നിയന്ത്രണം}, control.

\begin{cousinweb}
\textbf{\ml{നിയന്ത്രിക്കുക}} (\emph{niyantrikkuka}) is built from Sanskrit \emph{niyantraṇa}, \textquotedblleft{}restraining, checking, governing\textquotedblright{} (Monier-Williams), with the Malayalam verb ending -\ml{ിക്കുക}.
\end{cousinweb}

\begin{grammarlens}[title={What you've built}]
One more word for saying what happens if.
\end{grammarlens}

\subsection*{Your turn}
\begin{itemize}
\raggedright
  \item \emph{Say it:} \emph{niyantrikkuka}
  \item \emph{Say it:} \emph{niyantrikkuka}, once more
  \item \emph{Say it:} \emph{pōlīsŭ āḷkkūṭṭatte niyantriccu.} — The police controlled the crowd.
  \item \emph{Recall:} say the Malayalam for to analyse, then the Malayalam for a method, then say \emph{niyantrikkuka} again
\end{itemize}

\begin{tcolorbox}[breakable,colback=teal!4,colframe=teal!35!black,title={Before you move on}]
What does \ml{നിയന്ത്രിക്കുക} mean? (\textquotedblleft{}To control\textquotedblright{}.) Say it once more.
\end{tcolorbox}
\section[\texorpdfstring{\ml{വർദ്ധിക്കുക} (varddhikkuka)}{വർദ്ധിക്കുക (varddhikkuka)}]{\ml{വർദ്ധിക്കുക} (varddhikkuka) — to go up; to grow}

\label{lesson:ML-C449-varddhikkuka}



\begin{quote}
Before the new one: say the Malayalam for a method, then the Malayalam for to control.
\end{quote}

\subsection*{You'll want to know: \ml{വർദ്ധിക്കുക}}
\textbf{\ml{വർദ്ധിക്കുക}} — \emph{varddhikkuka} — \textquotedblleft{}to go up; to grow\textquotedblright{}.

\begin{quote}
\ml{നഗരത്തിൽ} \ml{ചൂട്} \ml{വർഷം} \ml{തോറും} \ml{വർദ്ധിക്കുന്നു}.
\end{quote}

\begin{quote}
\emph{nagarattil cūṭŭ varṣaṁ tōṟuṁ varddhikkunnu.}
\end{quote}

\begin{quote}
The heat in the city goes up year by year.
\end{quote}

Past: \ml{വർദ്ധിച്ചു}, \emph{varddhiccu}.

In talk people often say \ml{കൂടുക}.

\begin{cousinweb}
\textbf{\ml{വർദ്ധിക്കുക}} (\emph{varddhikkuka}) is built from the Sanskrit root \emph{vṛdh}, \textquotedblleft{}to grow\textquotedblright{}, as in \emph{vardhayati}, \textquotedblleft{}makes grow\textquotedblright{} (CDIAL 11382), with the Malayalam verb ending -\ml{ിക്കുക}.
\end{cousinweb}

\begin{grammarlens}[title={What you've built}]
One more word for saying what happens if.
\end{grammarlens}

\subsection*{Your turn}
\begin{itemize}
\raggedright
  \item \emph{Say it:} \emph{varddhikkuka}
  \item \emph{Say it:} \emph{varddhikkuka}, once more
  \item \emph{Say it:} \emph{nagarattil cūṭŭ varṣaṁ tōṟuṁ varddhikkunnu.} — The heat in the city goes up year by year.
  \item \emph{Recall:} say the Malayalam for a method, then the Malayalam for to control, then say \emph{varddhikkuka} again
\end{itemize}

\begin{tcolorbox}[breakable,colback=teal!4,colframe=teal!35!black,title={Before you move on}]
What does \ml{വർദ്ധിക്കുക} mean? (\textquotedblleft{}To go up; to grow\textquotedblright{}.) Say it once more.
\end{tcolorbox}
\section[\texorpdfstring{\ml{വർദ്ധിപ്പിക്കുക} (varddhippikkuka)}{വർദ്ധിപ്പിക്കുക (varddhippikkuka)}]{\ml{വർദ്ധിപ്പിക്കുക} (varddhippikkuka) — to increase; to raise}

\label{lesson:ML-C449-varddhippikkuka}



\begin{quote}
Before the new one: say the Malayalam for to control, then the Malayalam for to go up.
\end{quote}

\subsection*{You'll want to know: \ml{വർദ്ധിപ്പിക്കുക}}
\textbf{\ml{വർദ്ധിപ്പിക്കുക}} — \emph{varddhippikkuka} — \textquotedblleft{}to increase; to raise\textquotedblright{}.

\begin{quote}
\ml{കമ്പനി} \ml{ശമ്പളം} \ml{പത്ത്} \ml{ശതമാനം} \ml{വർദ്ധിപ്പിച്ചു}.
\end{quote}

\begin{quote}
\emph{kampani śampaḷaṁ pattŭ śatamānaṁ varddhippiccu.}
\end{quote}

\begin{quote}
The company raised salaries by ten per cent.
\end{quote}

Past: \ml{വർദ്ധിപ്പിച്ചു}, \emph{varddhippiccu}.

The causative of \ml{വർദ്ധിക്കുക}: to make a thing grow.

\begin{cousinweb}
\textbf{\ml{വർദ്ധിപ്പിക്കുക}} (\emph{varddhippikkuka}) is built from Sanskrit \emph{vardhayati}, \textquotedblleft{}makes grow, increases\textquotedblright{} (CDIAL 11382), with the Malayalam causative verb ending -\ml{ിപ്പിക്കുക}.
\end{cousinweb}

\begin{grammarlens}[title={What you've built}]
One more word for saying what happens if.
\end{grammarlens}

\subsection*{Your turn}
\begin{itemize}
\raggedright
  \item \emph{Say it:} \emph{varddhippikkuka}
  \item \emph{Say it:} \emph{varddhippikkuka}, once more
  \item \emph{Say it:} \emph{kampani śampaḷaṁ pattŭ śatamānaṁ varddhippiccu.} — The company raised salaries by ten per cent.
  \item \emph{Recall:} say the Malayalam for to control, then the Malayalam for to go up, then say \emph{varddhippikkuka} again
\end{itemize}

\begin{tcolorbox}[breakable,colback=teal!4,colframe=teal!35!black,title={Before you move on}]
What does \ml{വർദ്ധിപ്പിക്കുക} mean? (\textquotedblleft{}To increase; to raise\textquotedblright{}.) Say it once more.
\end{tcolorbox}
\section[\texorpdfstring{\ml{ഉയരുക} (uyaruka)}{ഉയരുക (uyaruka)}]{\ml{ഉയരുക} (uyaruka) — to go higher; to rise}

\label{lesson:ML-C449-uyaruka}



\begin{quote}
Before the new one: say the Malayalam for to go up, then the Malayalam for to increase.
\end{quote}

\subsection*{You'll want to know: \ml{ഉയരുക}}
\textbf{\ml{ഉയരുക}} — \emph{uyaruka} — \textquotedblleft{}to go higher; to rise\textquotedblright{}.

\begin{quote}
\ml{പുഴയിലെ} \ml{വെള്ളം} \ml{രാത്രി} \ml{ഉയർന്നു}.
\end{quote}

\begin{quote}
\emph{puḻayile veḷḷaṁ rātri uyarnnu.}
\end{quote}

\begin{quote}
The water in the river rose in the night.
\end{quote}

Past: \ml{ഉയർന്നു}, \emph{uyarnnu}.

Its past is \ml{ഉയർന്നു}. The noun \ml{ഉയരം}, height, is the same root.

\begin{cousinweb}
\textbf{\ml{ഉയരുക}} (\emph{uyaruka}) is native Dravidian: DEDR 646 lists Malayalam \emph{uyaruka}, \textquotedblleft{}to rise, to be high\textquotedblright{}, with \emph{uyaram}, \textquotedblleft{}height\textquotedblright{}.
\end{cousinweb}

\begin{grammarlens}[title={What you've built}]
One more word for saying what happens if.
\end{grammarlens}

\subsection*{Your turn}
\begin{itemize}
\raggedright
  \item \emph{Say it:} \emph{uyaruka}
  \item \emph{Say it:} \emph{uyaruka}, once more
  \item \emph{Say it:} \emph{puḻayile veḷḷaṁ rātri uyarnnu.} — The water in the river rose in the night.
  \item \emph{Recall:} say the Malayalam for to go up, then the Malayalam for to increase, then say \emph{uyaruka} again
\end{itemize}

\begin{tcolorbox}[breakable,colback=teal!4,colframe=teal!35!black,title={Before you move on}]
What does \ml{ഉയരുക} mean? (\textquotedblleft{}To go higher; to rise\textquotedblright{}.) Say it once more.
\end{tcolorbox}
\section[\texorpdfstring{\ml{ഉയർത്തുക} (uyarttuka)}{ഉയർത്തുക (uyarttuka)}]{\ml{ഉയർത്തുക} (uyarttuka) — to lift up; to raise}

\label{lesson:ML-C449-uyarttuka}



\begin{quote}
Before the new one: say the Malayalam for to increase, then the Malayalam for to go higher.
\end{quote}

\subsection*{You'll want to know: \ml{ഉയർത്തുക}}
\textbf{\ml{ഉയർത്തുക}} — \emph{uyarttuka} — \textquotedblleft{}to lift up; to raise\textquotedblright{}.

\begin{quote}
\ml{ആ} \ml{ചോദ്യം} \ml{അവൾ} \ml{യോഗത്തിൽ} \ml{ഉയർത്തി}.
\end{quote}

\begin{quote}
\emph{ā cōdyaṁ avaḷ yōgattil uyartti.}
\end{quote}

\begin{quote}
She raised that question at the meeting.
\end{quote}

Past: \ml{ഉയർത്തി}, \emph{uyartti}.

The causative of \ml{ഉയരുക}: to make a thing rise.

\begin{cousinweb}
\textbf{\ml{ഉയർത്തുക}} (\emph{uyarttuka}) is native Dravidian: DEDR 646 lists Malayalam \emph{uyarttuka}, \textquotedblleft{}to raise, to elevate\textquotedblright{}, the causative of \ml{ഉയരുക}.
\end{cousinweb}

\begin{grammarlens}[title={What you've built}]
One more word for saying what happens if.
\end{grammarlens}

\subsection*{Your turn}
\begin{itemize}
\raggedright
  \item \emph{Say it:} \emph{uyarttuka}
  \item \emph{Say it:} \emph{uyarttuka}, once more
  \item \emph{Say it:} \emph{ā cōdyaṁ avaḷ yōgattil uyartti.} — She raised that question at the meeting.
  \item \emph{Recall:} say the Malayalam for to increase, then the Malayalam for to go higher, then say \emph{uyarttuka} again
\end{itemize}

\begin{tcolorbox}[breakable,colback=teal!4,colframe=teal!35!black,title={Before you move on}]
What does \ml{ഉയർത്തുക} mean? (\textquotedblleft{}To lift up; to raise\textquotedblright{}.) Say it once more.
\end{tcolorbox}
